% Options for packages loaded elsewhere
\PassOptionsToPackage{unicode}{hyperref}
\PassOptionsToPackage{hyphens}{url}
\PassOptionsToPackage{dvipsnames,svgnames,x11names}{xcolor}
\documentclass[
  french,
  11pt,
  a4paper,
]{article}
\usepackage{xcolor}
\usepackage[margin=2.2cm]{geometry}
\usepackage{amsmath,amssymb}
\setcounter{secnumdepth}{-\maxdimen} % remove section numbering
\usepackage{iftex}
\ifPDFTeX
  \usepackage[T1]{fontenc}
  \usepackage[utf8]{inputenc}
  \usepackage{textcomp} % provide euro and other symbols
\else % if luatex or xetex
  \usepackage{unicode-math} % this also loads fontspec
  \defaultfontfeatures{Scale=MatchLowercase}
  \defaultfontfeatures[\rmfamily]{Ligatures=TeX,Scale=1}
\fi
\usepackage{lmodern}
\ifPDFTeX\else
  % xetex/luatex font selection
\fi
% Use upquote if available, for straight quotes in verbatim environments
\IfFileExists{upquote.sty}{\usepackage{upquote}}{}
\IfFileExists{microtype.sty}{% use microtype if available
  \usepackage[]{microtype}
  \UseMicrotypeSet[protrusion]{basicmath} % disable protrusion for tt fonts
}{}
\makeatletter
\@ifundefined{KOMAClassName}{% if non-KOMA class
  \IfFileExists{parskip.sty}{%
    \usepackage{parskip}
  }{% else
    \setlength{\parindent}{0pt}
    \setlength{\parskip}{6pt plus 2pt minus 1pt}}
}{% if KOMA class
  \KOMAoptions{parskip=half}}
\makeatother
\usepackage{longtable,booktabs,array}
\usepackage{calc} % for calculating minipage widths
% Correct order of tables after \paragraph or \subparagraph
\usepackage{etoolbox}
\makeatletter
\patchcmd\longtable{\par}{\if@noskipsec\mbox{}\fi\par}{}{}
\makeatother
% Allow footnotes in longtable head/foot
\IfFileExists{footnotehyper.sty}{\usepackage{footnotehyper}}{\usepackage{footnote}}
\makesavenoteenv{longtable}
\ifLuaTeX
\usepackage[bidi=basic,shorthands=off,]{babel}
\else
\usepackage[bidi=default,shorthands=off,]{babel}
\fi
\ifLuaTeX
  \usepackage{selnolig} % disable illegal ligatures
\fi
\setlength{\emergencystretch}{3em} % prevent overfull lines
\providecommand{\tightlist}{%
  \setlength{\itemsep}{0pt}\setlength{\parskip}{0pt}}
% Fichier LaTeX généré par docs/build/build_pdf.sh à partir de 00_fiche_demo.md.
% Compilation (depuis le dossier de ce fichier), deux fois pour la table des matières :
%   pdflatex 05_Fiche_de_demonstration.tex
% Chaque image est un bloc \begin{figure} ... \includegraphics{fichier} ... \end{figure}.
% Pour mettre votre capture : changez seulement le nom du fichier dans \includegraphics{...}
% (chemin relatif à ce dossier, par exemple ../images/guides/ma_capture.png), puis recompilez.
% Une image absente n'est pas insérée : le texte reste, sans figure.
\newcommand{\CredixShortTitle}{Fiche de démonstration}
% Préambule commun aux PDF du dossier de remise CREDIX (compilé avec pdfLaTeX via pandoc)
\usepackage{helvet}
\renewcommand{\familydefault}{\sfdefault}
\usepackage{pifont}
\usepackage{amssymb}
% Caractères Unicode utilisés dans les documents (pdfLaTeX ne les connaît pas d'office)
\DeclareUnicodeCharacter{03C1}{\ensuremath{\rho}}
\DeclareUnicodeCharacter{2265}{\ensuremath{\geq}}
\DeclareUnicodeCharacter{2264}{\ensuremath{\leq}}
\DeclareUnicodeCharacter{2248}{\ensuremath{\approx}}
\DeclareUnicodeCharacter{2192}{\ensuremath{\rightarrow}}
\DeclareUnicodeCharacter{2190}{\ensuremath{\leftarrow}}
\DeclareUnicodeCharacter{2212}{\ensuremath{-}}
\DeclareUnicodeCharacter{2713}{\ding{51}}
\DeclareUnicodeCharacter{2714}{\ding{52}}
\DeclareUnicodeCharacter{2717}{\ding{55}}
\DeclareUnicodeCharacter{25B2}{\ensuremath{\blacktriangle}}
\DeclareUnicodeCharacter{25CB}{\ensuremath{\circ}}
\DeclareUnicodeCharacter{2500}{-}
\DeclareUnicodeCharacter{2502}{|}
\DeclareUnicodeCharacter{250C}{+}
\DeclareUnicodeCharacter{2510}{+}
\DeclareUnicodeCharacter{2514}{+}
\DeclareUnicodeCharacter{2518}{+}
\DeclareUnicodeCharacter{251C}{+}
\DeclareUnicodeCharacter{2524}{+}
\DeclareUnicodeCharacter{252C}{+}
\DeclareUnicodeCharacter{2534}{+}
\DeclareUnicodeCharacter{253C}{+}
\DeclareUnicodeCharacter{0192}{\textit{f}}
\DeclareUnicodeCharacter{2011}{-}
\DeclareUnicodeCharacter{202F}{\,}
\DeclareUnicodeCharacter{2009}{\,}
\DeclareUnicodeCharacter{00D7}{\ensuremath{\times}}
\DeclareUnicodeCharacter{25A0}{\ensuremath{\blacksquare}}
\DeclareUnicodeCharacter{2022}{\ensuremath{\bullet}}
\DeclareUnicodeCharacter{2026}{\ldots}
\DeclareUnicodeCharacter{2460}{\ding{172}}
\DeclareUnicodeCharacter{2461}{\ding{173}}
\DeclareUnicodeCharacter{2462}{\ding{174}}
\DeclareUnicodeCharacter{2463}{\ding{175}}
\DeclareUnicodeCharacter{2464}{\ding{176}}
\DeclareUnicodeCharacter{2465}{\ding{177}}
\DeclareUnicodeCharacter{2466}{\ding{178}}
\DeclareUnicodeCharacter{2467}{\ding{179}}
\usepackage{fvextra}
\DefineVerbatimEnvironment{verbatim}{Verbatim}{breaklines=true,breakanywhere=true,fontsize=\small,frame=single,framerule=0.3pt,rulecolor=\color[gray]{0.7},framesep=2mm,xleftmargin=1mm,xrightmargin=1mm}
\fvset{breaklines=true,breakanywhere=true}
\usepackage{fancyhdr}
\usepackage{tcolorbox}
\usepackage{colortbl}
\usepackage{enumitem}
\definecolor{credixblue}{HTML}{1F3A5F}
\definecolor{credixgray}{HTML}{5F6B7A}
\definecolor{boxbg}{HTML}{F3F5F8}
\definecolor{credixaccent}{HTML}{D9822B}
\usepackage{tikz}
\setlength{\emergencystretch}{4em}
\renewcommand{\arraystretch}{1.3}
\setlength{\parindent}{0pt}
\setlength{\parskip}{0.55em}
\setlist{itemsep=0.15em,topsep=0.3em}
\pagestyle{fancy}
\fancyhf{}
\fancyhead[L]{\small\color{credixgray}\textsc{\CredixShortTitle}}
\fancyhead[R]{\small\color{credixgray}Dossier de remise}
\fancyfoot[C]{\small\color{credixgray}\thepage}
\renewcommand{\headrulewidth}{0.3pt}
\renewcommand{\footrulewidth}{0pt}
\fancypagestyle{plain}{\fancyhf{}\fancyfoot[C]{\small\color{credixgray}\thepage}\renewcommand{\headrulewidth}{0pt}}
\usepackage{titlesec}
\titleformat{\section}{\Large\bfseries\color{credixblue}}{}{0em}{}[\vspace{-0.2em}\color{credixblue}\titlerule]
\titleformat{\subsection}{\large\bfseries\color{credixblue}}{}{0em}{}
\titleformat{\subsubsection}{\normalsize\bfseries\color{credixblue}}{}{0em}{}
% Cadre de repère pour une capture manquante
\newcommand{\CaptureManquante}[2]{%
  \begin{tcolorbox}[colback=boxbg,colframe=credixgray,boxrule=0.5pt,arc=1mm,left=3mm,right=3mm,top=3mm,bottom=3mm,halign=center]
  \textbf{CAPTURE À INSÉRER}\\[0.4em]#1\\[0.4em]{\footnotesize\ttfamily #2}
  \end{tcolorbox}}

% Figures : chaque image est un bloc \begin{figure} ... \includegraphics{fichier} ... \end{figure}.
% Si le fichier est absent, un cadre « CAPTURE À INSÉRER » s'affiche à sa place et la compilation continue.
% Pour mettre votre image : changez seulement le nom du fichier dans \includegraphics{...}.
\usepackage{graphicx}
\usepackage{float}
\usepackage{needspace}
\newcommand{\CredixAlt}{Capture}
\let\CredixOrigIncludegraphics\includegraphics
\renewcommand{\includegraphics}[2][]{%
  \IfFileExists{#2}%
    {\CredixOrigIncludegraphics[#1]{#2}}%
    {\CaptureManquante{\CredixAlt}{\detokenize{#2}}}}
\usepackage{bookmark}
\IfFileExists{xurl.sty}{\usepackage{xurl}}{} % add URL line breaks if available
\urlstyle{same}
\hypersetup{
  pdftitle={Fiche de démonstration},
  pdfauthor={SINGHE PENKA Hendrix Donavan},
  pdflang={fr},
  colorlinks=true,
  linkcolor={credixblue},
  filecolor={Maroon},
  citecolor={Blue},
  urlcolor={credixblue},
  pdfcreator={LaTeX via pandoc}}

\author{}
\date{}

\begin{document}

\begin{titlepage}
\thispagestyle{empty}
\begin{tikzpicture}[remember picture,overlay]
  \fill[credixblue] (current page.north west) rectangle ([yshift=-8cm]current page.north east);
  \fill[credixaccent] ([yshift=-8cm]current page.north west) rectangle ([yshift=-8.25cm]current page.north east);
\end{tikzpicture}
\noindent
{\color{white}
\vspace*{0.3cm}
{\large\bfseries\textsc{Projet de fin d'études}\par}
\vspace{0.6cm}
{\Large École Nationale Supérieure Polytechnique de Yaoundé\par}
\vspace{0.15cm}
{\Large Université de Yaoundé I\par}
\vspace{0.9cm}
{\normalsize Diplôme d'Ingénieur de Conception en Génie Informatique\par}
{\normalsize Année académique 2025--2026\par}
}
\vspace{2.8cm}
\noindent{\fontsize{28}{34}\selectfont\bfseries\color{credixblue} Fiche de démonstration\par}
\vspace{0.6cm}
\noindent{\Large\color{credixgray} Déroulé de la séance avec l'encadrant\par}
\vspace{0.8cm}
\noindent{\color{credixaccent}\rule{4cm}{2.5pt}\par}
\vfill
\noindent{\small\color{credixgray} Présenté par\par}
\vspace{0.1cm}
\noindent{\Large\bfseries SINGHE PENKA Hendrix Donavan\par}
\vspace{0.7cm}
\noindent{\small\color{credixgray} Version du 19/09/2026\par}
\end{titlepage}


\renewcommand*\contentsname{Table des matières}
{
\setcounter{tocdepth}{2}
\tableofcontents
}
\section{Fiche de démonstration}\label{fiche-de-duxe9monstration}

\textbf{Déroulé de la séance avec l\textquotesingle encadrant : vérifier
le projet, comprendre la méthode, tester la méthode sur un exemple.}

\begin{center}\rule{0.5\linewidth}{0.5pt}\end{center}

\subsection{1. Objectifs et déroulé de la
séance}\label{1-objectifs-et-duxe9rouluxe9-de-la-suxe9ance}

L\textquotesingle encadrant veut vérifier \textbf{trois choses}, dans
cet ordre :

\begin{longtable}[]{@{}
  >{\raggedright\arraybackslash}p{(\linewidth - 6\tabcolsep) * \real{0.1430}}
  >{\raggedright\arraybackslash}p{(\linewidth - 6\tabcolsep) * \real{0.3289}}
  >{\raggedright\arraybackslash}p{(\linewidth - 6\tabcolsep) * \real{0.3289}}
  >{\raggedright\arraybackslash}p{(\linewidth - 6\tabcolsep) * \real{0.1692}}@{}}
\toprule\noalign{}
\begin{minipage}[b]{\linewidth}\raggedright
Partie
\end{minipage} & \begin{minipage}[b]{\linewidth}\raggedright
Question de l\textquotesingle encadrant
\end{minipage} & \begin{minipage}[b]{\linewidth}\raggedright
Ce qu\textquotesingle on lui montre
\end{minipage} & \begin{minipage}[b]{\linewidth}\raggedright
Document
\end{minipage} \\
\midrule\noalign{}
\endhead
\bottomrule\noalign{}
\endlastfoot
\textbf{1} & \emph{« Ce qui a été fait est-il bien sur Git,
fonctionne-t-il et est-il utilisable ? »} & Le projet récupéré depuis
GitHub, lancé et testé & README \\
\textbf{2} & \emph{« Comment avez-vous procédé pour y arriver ? »} & La
méthode de travail avec Claude, avec les vrais documents du projet &
Guide 2 (et guide 1 pour les outils) \\
\textbf{3} & \emph{« Cette méthode marche-t-elle vraiment ?
Montrez-le-moi. »} & Une conception réduite (MVP) confiée à Claude Code,
en direct & Guide 3 (MVP) \\
\end{longtable}

\textbf{Répartition indicative du temps} (à ajuster) : partie 1, environ
15 minutes ; partie 2, environ 25 minutes ; partie 3, environ 40
minutes.

\textbf{Ordre conseillé.} L\textquotesingle installation de la partie 1
est \textbf{longue} (jusqu\textquotesingle à 50 minutes la première
fois). On ne l\textquotesingle attend donc pas sans rien faire :

\begin{itemize}
\tightlist
\item
  \textbf{Si l\textquotesingle encadrant a déjà installé le projet} :
  parties 1, 2, puis 3.
\item
  \textbf{Sinon} : dès les premières minutes, il lance la commande
  d\textquotesingle installation du README (section 5.2). Pendant
  qu\textquotesingle elle travaille, vous faites la \textbf{partie 2},
  puis la \textbf{partie 3}. Quand l\textquotesingle installation est
  terminée, vous faites la \textbf{partie 1} (les tests).
\end{itemize}

\begin{center}\rule{0.5\linewidth}{0.5pt}\end{center}

\subsection{2. Avant la séance : les préparatifs
indispensables}\label{2-avant-la-suxe9ance--les-pruxe9paratifs-indispensables}

\begin{quote}
\textbf{Le point critique : la première installation est longue.} La
construction des conteneurs Docker télécharge environ 700 Mo et a duré
\textbf{environ 50 minutes} sur une connexion lente lors de nos essais.
Elle ne doit \textbf{jamais} être lancée en direct pendant la séance.
\textbf{Elle doit être faite avant.}
\end{quote}

\subsubsection{2.1 Accès au dépôt (à faire la
veille)}\label{21-accuxe8s-au-duxe9puxf4t-uxe0-faire-la-veille}

\begin{enumerate}
\def\labelenumi{\arabic{enumi}.}
\tightlist
\item
  \textbf{Inviter l\textquotesingle encadrant} sur GitHub : ouvrir le
  dépôt, puis \emph{Settings}, \emph{Collaborators}, \emph{Add people},
  et saisir son nom d\textquotesingle utilisateur GitHub. Le dépôt est
  \textbf{privé} : sans invitation acceptée, il ne peut pas le cloner.
\item
  Lui demander de \textbf{créer un jeton d\textquotesingle accès
  personnel} (README, section 4.1) : sans lui, Git refuse le mot de
  passe du compte.
\item
  Se rappeler qu\textquotesingle un collaborateur d\textquotesingle un
  dépôt personnel peut \textbf{lire et modifier} le dépôt.
\end{enumerate}

\subsubsection{2.2 Ordinateur de l\textquotesingle encadrant (idéalement
la veille)}\label{22-ordinateur-de-lencadrant-iduxe9alement-la-veille}

\begin{longtable}[]{@{}
  >{\raggedright\arraybackslash}p{(\linewidth - 2\tabcolsep) * \real{0.2071}}
  >{\raggedright\arraybackslash}p{(\linewidth - 2\tabcolsep) * \real{0.7629}}@{}}
\toprule\noalign{}
\begin{minipage}[b]{\linewidth}\raggedright
À vérifier
\end{minipage} & \begin{minipage}[b]{\linewidth}\raggedright
Minimum
\end{minipage} \\
\midrule\noalign{}
\endhead
\bottomrule\noalign{}
\endlastfoot
Système & Linux 64 bits \\
Espace disque libre & 10 Go (la construction a besoin
d\textquotesingle environ 4 Go de marge temporaire à la fin) \\
Outils & Git, Docker et Docker Compose, Node.js 18.18 ou plus, curl,
OpenSSL (README, section 3) \\
Ports libres & 3001, 8080, 8001, 5000, 27017 \\
\end{longtable}

\textbf{Lancer, la veille, la première installation} en suivant le
README (sections 4 à 6). Le README indique, à chaque étape, ce
qu\textquotesingle il faut voir. Le lancement suivant ne prend
qu\textquotesingle une minute environ.

\subsubsection{2.3 Pour la partie 3 (Claude
Code)}\label{23-pour-la-partie-3-claude-code}

\begin{itemize}
\tightlist
\item
  \textbf{VS Code 1.94 ou plus} avec l\textquotesingle extension
  \textbf{Claude Code} installée (guide 1, section 7) ;
\item
  un \textbf{abonnement Claude payant} (Pro au minimum) ou un compte
  Console : sans lui, Claude Code ne démarre pas ;
\item
  \textbf{être connecté} dans l\textquotesingle extension avant la
  séance ;
\item
  avoir \textbf{copié}
  \texttt{docs/\allowbreak{}guides/\allowbreak{}03\_\allowbreak{}conception\_\allowbreak{}MVP.\allowbreak{}md}
  dans un dossier de travail vide, sous le nom
  \texttt{docs/\allowbreak{}conception\_\allowbreak{}mvp.\allowbreak{}md}.
\end{itemize}

\subsubsection{2.4 Côté auteur}\label{24-cuxf4tuxe9-auteur}

\begin{itemize}
\tightlist
\item
  Garder \textbf{votre propre pile en marche} (Docker et frontend) :
  elle sert de \textbf{plan B} si l\textquotesingle installation de
  l\textquotesingle encadrant n\textquotesingle est pas terminée.
\item
  Avoir sous la main : le lien du dépôt, le README en PDF, les guides en
  PDF, le journal et les documents de passation du projet,
  l\textquotesingle abonnement Claude.
\item
  Une connexion Internet stable.
\end{itemize}

\subsubsection{2.5 Plan B}\label{25-plan-b}

\begin{longtable}[]{@{}
  >{\raggedright\arraybackslash}p{(\linewidth - 2\tabcolsep) * \real{0.4354}}
  >{\raggedright\arraybackslash}p{(\linewidth - 2\tabcolsep) * \real{0.5346}}@{}}
\toprule\noalign{}
\begin{minipage}[b]{\linewidth}\raggedright
Problème
\end{minipage} & \begin{minipage}[b]{\linewidth}\raggedright
Solution
\end{minipage} \\
\midrule\noalign{}
\endhead
\bottomrule\noalign{}
\endlastfoot
L\textquotesingle installation de l\textquotesingle encadrant
n\textquotesingle est pas terminée & Faire la partie 1 \textbf{sur votre
machine} (partage d\textquotesingle écran), et laisser
l\textquotesingle encadrant terminer son installation pendant les
parties 2 et 3 \\
Pas d\textquotesingle abonnement Claude disponible pour
l\textquotesingle encadrant & Faire la partie 3 \textbf{sur votre
compte} (partage d\textquotesingle écran), l\textquotesingle encadrant
relit et dirige \\
Claude Code est lent ou saturé (limite d\textquotesingle usage atteinte)
& Vérifier votre consommation (guide 1, section 4.3) ; se limiter au
plan et au module M2 \\
\end{longtable}

\begin{center}\rule{0.5\linewidth}{0.5pt}\end{center}

\subsection{3. Partie 1 : vérifier que le projet fonctionne depuis
Git}\label{3-partie-1--vuxe9rifier-que-le-projet-fonctionne-depuis-git}

\textbf{Objectif :} l\textquotesingle encadrant constate que ce qui est
sur GitHub se lance et se teste.

\subsubsection{3.1 Déroulé}\label{31-duxe9rouluxe9}

\begin{longtable}[]{@{}
  >{\raggedright\arraybackslash}p{(\linewidth - 4\tabcolsep) * \real{0.1430}}
  >{\raggedright\arraybackslash}p{(\linewidth - 4\tabcolsep) * \real{0.6065}}
  >{\raggedright\arraybackslash}p{(\linewidth - 4\tabcolsep) * \real{0.2205}}@{}}
\toprule\noalign{}
\begin{minipage}[b]{\linewidth}\raggedright
Étape
\end{minipage} & \begin{minipage}[b]{\linewidth}\raggedright
Action
\end{minipage} & \begin{minipage}[b]{\linewidth}\raggedright
Référence
\end{minipage} \\
\midrule\noalign{}
\endhead
\bottomrule\noalign{}
\endlastfoot
1 & Cloner le projet & README, section 4.2 \\
2 & Vérifier la copie : 16 fichiers de modèle & README, section 4.3 \\
3 & Vérifier l\textquotesingle état du backend :
\texttt{docker\ compose\ ps}, les logs, \texttt{/health} & README,
section 5.3 \\
4 & Premier lancement (administrateur et données de démonstration) &
README, section 5.4 \\
5 & Vérifier le frontend & README, section 6.4 \\
6 & Parcours de test : connexion, création d\textquotesingle un agent,
scoring de \texttt{HC-100001}, dossier en revue avec \texttt{HC-100241},
journal d\textquotesingle audit & README, section 7 \\
\end{longtable}

\textbf{Comptes de test.} \texttt{admin}, \texttt{agent.test} et
\texttt{superviseur.\allowbreak{}test}, tous les trois avec le mot de
passe d\textquotesingle exemple \texttt{Demo12345} (README, sections 5.4
et 7). Ce mot de passe est sans risque : il ne sert que sur la machine
de test.

\subsubsection{3.2 Ce que l\textquotesingle encadrant doit
constater}\label{32-ce-que-lencadrant-doit-constater}

\begin{longtable}[]{@{}
  >{\raggedright\arraybackslash}p{(\linewidth - 2\tabcolsep) * \real{0.2670}}
  >{\raggedright\arraybackslash}p{(\linewidth - 2\tabcolsep) * \real{0.7030}}@{}}
\toprule\noalign{}
\begin{minipage}[b]{\linewidth}\raggedright
Contrôle
\end{minipage} & \begin{minipage}[b]{\linewidth}\raggedright
Réussi si\ldots{}
\end{minipage} \\
\midrule\noalign{}
\endhead
\bottomrule\noalign{}
\endlastfoot
Récupération & Le \texttt{git\ clone} réussit ; `ls
scoring-backend/artefacts \\
Backend & 4 conteneurs \texttt{Up}, l\textquotesingle API
\texttt{healthy} ; \texttt{/health} : \texttt{"statut":"ok"}, 27
variables, 27 phrases, \texttt{"mock\_\allowbreak{}mode":false} \\
Frontend & La page de connexion s\textquotesingle affiche ; les trois
espaces (Agent, Superviseur, Administrateur) sont accessibles \\
Scoring & \texttt{HC-100001} : score 603, \texttt{ACCORDÉ} ;
\texttt{HC-100241} : score 571, \texttt{EN\ REVUE} ; \texttt{HC-101327}
: score 536, \texttt{REFUSÉ} \\
Traçabilité & Le journal d\textquotesingle audit montre connexions,
création de compte, scorings et décision du superviseur \\
\end{longtable}

\subsubsection{3.3 Si quelque chose ne marche
pas}\label{33-si-quelque-chose-ne-marche-pas}

Ne pas improviser : ouvrir la \textbf{section 10 du README} (dépannage),
qui liste les symptômes et les solutions vérifiées. Les cas les plus
probables sont un port déjà utilisé, un manque d\textquotesingle espace
disque, et l\textquotesingle oubli de redémarrer l\textquotesingle API
après le chargement des données.

\begin{center}\rule{0.5\linewidth}{0.5pt}\end{center}

\subsection{4. Partie 2 : comprendre comment le projet a été
conçu}\label{4-partie-2--comprendre-comment-le-projet-a-uxe9tuxe9-conuxe7u}

\textbf{Objectif :} l\textquotesingle encadrant comprend la
\textbf{méthode}, pas seulement le résultat. On s\textquotesingle appuie
sur le \textbf{guide 2}.

\subsubsection{4.1 Les cinq messages à faire
passer}\label{41-les-cinq-messages-uxe0-faire-passer}

\begin{longtable}[]{@{}
  >{\raggedright\arraybackslash}p{(\linewidth - 4\tabcolsep) * \real{0.0965}}
  >{\raggedright\arraybackslash}p{(\linewidth - 4\tabcolsep) * \real{0.4367}}
  >{\raggedright\arraybackslash}p{(\linewidth - 4\tabcolsep) * \real{0.4367}}@{}}
\toprule\noalign{}
\begin{minipage}[b]{\linewidth}\raggedright
N°
\end{minipage} & \begin{minipage}[b]{\linewidth}\raggedright
Message
\end{minipage} & \begin{minipage}[b]{\linewidth}\raggedright
Preuve à montrer
\end{minipage} \\
\midrule\noalign{}
\endhead
\bottomrule\noalign{}
\endlastfoot
\textbf{1} & On \textbf{conçoit avant de coder}, et on écrit tout en
Markdown & Le cahier des charges et la conception ; le README
lui-même \\
\textbf{2} & \textbf{Claude propose, l\textquotesingle humain décide} :
aucun livrable sans accord explicite & Les instructions du projet
(section 5, workflow de collaboration) \\
\textbf{3} & Le projet tient sur des \textbf{sessions}, grâce à des
documents : instructions, journal, passations, registre des décisions &
Le prompt de démarrage de session, le rapport de passation, le
journal \\
\textbf{4} & Chaque \textbf{choix} est justifié et classé par
\textbf{niveau de preuve} & L\textquotesingle exemple du WoE (guide 2,
section 8) \\
\textbf{5} & Claude \textbf{se trompe} ; la méthode permet de \textbf{le
voir} & Le tableau des erreurs réelles (guide 2, section 14) \\
\end{longtable}

\subsubsection{4.2 Les fichiers à ouvrir pendant la
présentation}\label{42-les-fichiers-uxe0-ouvrir-pendant-la-pruxe9sentation}

\begin{itemize}
\tightlist
\item
  les \textbf{instructions du projet} Claude (six sections) ;
\item
  un \textbf{prompt de démarrage de session} et un \textbf{rapport de
  passation} ;
\item
  le \textbf{journal de bord} et le \textbf{journal de rédaction}
  (règles actées, pièges) ;
\item
  le \textbf{registre des décisions méthodologiques} ;
\item
  le fichier \texttt{scanner.sh} et l\textquotesingle instantané de
  contexte qu\textquotesingle il produit.
\end{itemize}

\subsubsection{4.3 La démonstration des
outils}\label{43-la-duxe9monstration-des-outils}

Montrer, dans l\textquotesingle ordre : le \textbf{Projet} claude.ai
(instructions, base de connaissances) ; la page \textbf{Utilisation}
(limites de session) ; \textbf{Claude Code} dans VS Code et son
\textbf{mode Plan} (guide 1).

\begin{center}\rule{0.5\linewidth}{0.5pt}\end{center}

\subsection{5. Partie 3 : tester la méthode sur le
MVP}\label{5-partie-3--tester-la-muxe9thode-sur-le-mvp}

\textbf{Objectif :} montrer que la méthode \textbf{fonctionne sur un cas
qu\textquotesingle on ne connaît pas}. L\textquotesingle encadrant
regarde, ou fait lui-même, ce qui suit. On s\textquotesingle appuie sur
le \textbf{guide 3}.

\subsubsection{5.1 Protocole pas à pas}\label{51-protocole-pas-uxe0-pas}

\begin{longtable}[]{@{}
  >{\raggedright\arraybackslash}p{(\linewidth - 4\tabcolsep) * \real{0.1430}}
  >{\raggedright\arraybackslash}p{(\linewidth - 4\tabcolsep) * \real{0.6220}}
  >{\raggedright\arraybackslash}p{(\linewidth - 4\tabcolsep) * \real{0.2050}}@{}}
\toprule\noalign{}
\begin{minipage}[b]{\linewidth}\raggedright
Étape
\end{minipage} & \begin{minipage}[b]{\linewidth}\raggedright
Action
\end{minipage} & \begin{minipage}[b]{\linewidth}\raggedright
Durée indicative
\end{minipage} \\
\midrule\noalign{}
\endhead
\bottomrule\noalign{}
\endlastfoot
\textbf{1} & Créer un dossier vide \texttt{demo\_mvp},
l\textquotesingle ouvrir dans VS Code & 1 min \\
\textbf{2} & Y copier
\texttt{03\_\allowbreak{}conception\_\allowbreak{}MVP.\allowbreak{}md}
sous
\texttt{docs/\allowbreak{}conception\_\allowbreak{}mvp.\allowbreak{}md}
& 1 min \\
\textbf{3} & Ouvrir Claude Code, \textbf{passer en mode Plan}
(indicateur de mode) & 1 min \\
\textbf{4} & Donner le \textbf{prompt de l\textquotesingle annexe B} du
guide 3 & 1 min \\
\textbf{5} & \textbf{Lire le plan} proposé et le juger avec la grille
ci-dessous & 5 à 10 min \\
\textbf{6} & Corriger si besoin (commentaires dans le plan), puis
\textbf{approuver} & 5 min \\
\textbf{7} & Demander la réalisation du \textbf{module M2} (données et
modèle) & 5 à 10 min \\
\textbf{8} & Vérifier le résultat : l\textquotesingle AUC de test et la
reproductibilité & 3 min \\
\end{longtable}

\subsubsection{5.2 Grille d\textquotesingle évaluation du plan (étape
5)}\label{52-grille-duxe9valuation-du-plan-uxe9tape-5}

\begin{longtable}[]{@{}
  >{\raggedright\arraybackslash}p{(\linewidth - 2\tabcolsep) * \real{0.8595}}
  >{\raggedright\arraybackslash}p{(\linewidth - 2\tabcolsep) * \real{0.1105}}@{}}
\toprule\noalign{}
\begin{minipage}[b]{\linewidth}\raggedright
Question
\end{minipage} & \begin{minipage}[b]{\linewidth}\raggedright
Oui / Non
\end{minipage} \\
\midrule\noalign{}
\endhead
\bottomrule\noalign{}
\endlastfoot
Le plan couvre-t-il les six modules (socle, données et modèle, service
de scoring, sécurité, API, interface) ? & \\
Les modules sont-ils \textbf{ordonnés} selon leurs dépendances (M1, puis
M2, puis M3\ldots) ? & \\
Chaque module a-t-il ses \textbf{fichiers}, ses \textbf{tests} et son
\textbf{critère de fin} ? & \\
Le plan \textbf{reprend-il les exigences} (EF et ENF) du document, sans
en inventer ? & \\
Signale-t-il des \textbf{ambiguïtés} du document ? (un bon signe) & \\
\textbf{Ne crée-t-il aucun fichier} à ce stade (mode Plan) ? & \\
\end{longtable}

\subsubsection{5.3 Critères de succès du module
M2}\label{53-crituxe8res-de-succuxe8s-du-module-m2}

\begin{longtable}[]{@{}
  >{\raggedright\arraybackslash}p{(\linewidth - 2\tabcolsep) * \real{0.4261}}
  >{\raggedright\arraybackslash}p{(\linewidth - 2\tabcolsep) * \real{0.5439}}@{}}
\toprule\noalign{}
\begin{minipage}[b]{\linewidth}\raggedright
Vérification
\end{minipage} & \begin{minipage}[b]{\linewidth}\raggedright
Attendu
\end{minipage} \\
\midrule\noalign{}
\endhead
\bottomrule\noalign{}
\endlastfoot
Le script de génération produit 3 000 clients & Fichier
\texttt{clients\_\allowbreak{}synthetiques.\allowbreak{}csv} de 3 000
lignes \\
L\textquotesingle AUC de test affichée & \textbf{Au moins 0,70} (environ
0,75 avec les données de l\textquotesingle annexe A) \\
Deux exécutions successives & Mêmes coefficients à la troisième
décimale \\
Les tests du module & Ils passent \\
\end{longtable}

\subsubsection{5.4 Si le plan ou le module n\textquotesingle est pas
satisfaisant}\label{54-si-le-plan-ou-le-module-nest-pas-satisfaisant}

C\textquotesingle est \textbf{normal et instructif} :
c\textquotesingle est ce que la méthode est censée révéler.

\begin{enumerate}
\def\labelenumi{\arabic{enumi}.}
\tightlist
\item
  \textbf{Identifier ce qui manque} dans le plan (un module oublié, un
  ordre incohérent).
\item
  \textbf{Le dire à Claude} en mode Plan (« il manque \ldots{} », « M3
  dépend de M2 »), ou par commentaire dans le plan.
\item
  Si le document lui-même est ambigu, \textbf{le corriger} : la
  conception est la source de vérité.
\item
  \textbf{Recommencer le plan}, puis approuver.
\end{enumerate}

\subsubsection{5.5 Ce que la séance ne cherche pas à
prouver}\label{55-ce-que-la-suxe9ance-ne-cherche-pas-uxe0-prouver}

Que Claude écrit un système parfait en une passe. Elle cherche à prouver
que \textbf{la méthode (conception écrite, mode Plan, un module à la
fois, tests)} donne un travail \textbf{vérifiable et reprenable}.

\begin{center}\rule{0.5\linewidth}{0.5pt}\end{center}

\subsection{6. Questions probables et éléments de
réponse}\label{6-questions-probables-et-uxe9luxe9ments-de-ruxe9ponse}

\begin{longtable}[]{@{}
  >{\raggedright\arraybackslash}p{(\linewidth - 2\tabcolsep) * \real{0.4134}}
  >{\raggedright\arraybackslash}p{(\linewidth - 2\tabcolsep) * \real{0.5566}}@{}}
\toprule\noalign{}
\begin{minipage}[b]{\linewidth}\raggedright
Question de l\textquotesingle encadrant
\end{minipage} & \begin{minipage}[b]{\linewidth}\raggedright
Éléments de réponse
\end{minipage} \\
\midrule\noalign{}
\endhead
\bottomrule\noalign{}
\endlastfoot
\emph{Pourquoi concevoir avant de coder ?} & Une conception écrite est
relisible et corrigeable à coût presque nul ; du code écrit sans
conception cache ses erreurs. La conception est aussi la \textbf{source
de vérité} que Claude Code lit. \\
\emph{Comment savez-vous que Claude ne s\textquotesingle est pas trompé
?} & On \textbf{ne le croit pas sur parole} : chiffres vérifiés contre
l\textquotesingle exécution réelle, tests, relecture, registre des
décisions. Des erreurs réelles ont été trouvées ainsi (guide 2, section
14). \\
\emph{Que vaut la littérature apportée par Claude ?} & Elle est
\textbf{un point de départ}. Chaque référence importante est vérifiée,
et les incertaines restent marquées « à vérifier ». \\
\emph{Peut-on refaire la démarche sur un autre sujet ?} & Oui : le guide
2 donne les prompts, les modèles de documents et le cycle de session. Le
MVP en est un exemple complet. \\
\emph{Combien cela coûte-t-il ?} & Un abonnement Pro à 20 \$ par mois
suffit pour commencer ; Max (100 ou 200 \$) si l\textquotesingle usage
de Claude Code est intensif (guide 1, section 3). Prix vérifiés le 19
septembre 2026. \\
\emph{Et la confidentialité des données ?} & CREDIX a été conçu sur un
\textbf{jeu de données public} (Home Credit), sans données de clients
réels. Ne jamais confier de données réelles ni de secrets à un outil
externe sans autorisation. \\
\emph{Que fait CREDIX que le MVP ne fait pas ?} & Deux flux (score et
détection des profils atypiques), WoE, calibration, indice de couverture
ρc, trois rôles, base documentaire, PDF (guide 3, section 13). \\
\emph{Peut-on modifier le projet ?} & Oui : README, section 9 (Git), et
méthode du guide 2, section 13 (toujours en mode Plan). \\
\end{longtable}

\begin{center}\rule{0.5\linewidth}{0.5pt}\end{center}

\subsection{7. Liste de contrôle de la
séance}\label{7-liste-de-contruxf4le-de-la-suxe9ance}

\subsubsection{La veille}\label{la-veille}

\begin{itemize}
\tightlist
\item[$\square$]
  L\textquotesingle encadrant est \textbf{invité} sur le dépôt et a
  accepté l\textquotesingle invitation
\item[$\square$]
  Il a créé son \textbf{jeton d\textquotesingle accès}
\item[$\square$]
  Sa \textbf{première installation} est lancée, ou terminée
  (\texttt{docker\ compose\ ps} : 4 conteneurs, API \texttt{healthy})
\item[$\square$]
  VS Code, l\textquotesingle extension Claude Code et
  l\textquotesingle abonnement sont prêts
\item[$\square$]
  Votre propre pile fonctionne (plan B)
\item[$\square$]
  Espace disque suffisant (au moins 10 Go libres) des deux côtés
\end{itemize}

\subsubsection{Pendant la séance}\label{pendant-la-suxe9ance}

\begin{itemize}
\tightlist
\item[$\square$]
  Partie 1 : \texttt{git\ clone}, \texttt{/health}, connexion, un
  scoring
\item[$\square$]
  Partie 2 : les cinq messages, les vrais documents du projet
\item[$\square$]
  Partie 3 : plan en mode Plan, grille d\textquotesingle évaluation,
  module M2
\end{itemize}

\subsubsection{Après la séance}\label{apruxe8s-la-suxe9ance}

\begin{itemize}
\tightlist
\item[$\square$]
  Noter les remarques et corrections de l\textquotesingle encadrant
\item[$\square$]
  Mettre à jour les documents concernés (le \texttt{.md} est la source ;
  le \texttt{.tex} et le PDF se régénèrent avec
  \texttt{docs/\allowbreak{}build/\allowbreak{}build\_\allowbreak{}all.\allowbreak{}sh})
\item[$\square$]
  Enregistrer et envoyer les modifications avec Git (README, section 9)
\end{itemize}

\end{document}
